\subsection{مبانی نظری و فرمول‌بندی ریاضی}
فرض کنید ماتریس وزن اولیه لایه خطی با ابعاد خروجی $d_{\text{out}}$ و ابعاد ورودی $d_{\text{in}}$ به صورت $\mathbf{W} \in \mathbb{R}^{d_{\text{out}} \times d_{\text{in}}}$ داده شده باشد. تجزیه کامل مقادیر منفرد به صورت زیر است:
\begin{equation}
\mathbf{W} = \mathbf{U} \mathbf{\Sigma} \mathbf{V}^\top
\end{equation}
که در آن $\mathbf{U} \in \mathbb{R}^{d_{\text{out}} \times d_{\text{out}}}$ و $\mathbf{V} \in \mathbb{R}^{d_{\text{in}} \times d_{\text{in}}}$ ماتریس‌های متعامد بوده و $\mathbf{\Sigma}$ شامل مقادیر منفرد مرتب‌شده نزولی $\sigma_1 \ge \sigma_2 \ge \dots \ge \sigma_{\min(d_{\text{out}}, d_{\text{in}})} \ge 0$ است.

برای رتبه هدف $r \le \min(d_{\text{out}}, d_{\text{in}})$، بهترین تقریب رتبه $r$ حاصل می‌شود. به منظور حفظ توازن عددی در مقیاس گرادیان‌ها در گذر پسرو و پایداری لایه‌ها، انرژی ماتریس قطری $\mathbf{\Sigma}_r$ به طور متقارن بین فاکتور چپ و راست تقسیم می‌شود:
\begin{equation}
\mathbf{U}_{\text{factor}} = \mathbf{U}_{:, :r} \, \mathbf{\Sigma}_r^{1/2} \in \mathbb{R}^{d_{\text{out}} \times r}, \quad \mathbf{V}_{\text{factor}} = \mathbf{\Sigma}_r^{1/2} \, \mathbf{V}_{:, :r}^\top \in \mathbb{R}^{r \times d_{\text{in}}}
\end{equation}

بدین ترتیب، ضرب خطی لایه در زمان استنتاج به دو ضرب متوالی کم‌حجم تبدیل می‌گردد:
\begin{equation}
\mathbf{y} = (\mathbf{x} \mathbf{V}_{\text{factor}}^\top) \mathbf{U}_{\text{factor}}^\top + \mathbf{b}
\end{equation}

تعداد کل پارامترهای ذخیره‌شده از $d_{\text{in}} \cdot d_{\text{out}}$ به $r(d_{\text{in}} + d_{\text{out}})$ کاهش می‌یابد. شرط الزامی برای تحقق فشرده‌سازی پارامتری عبارت است از:
\begin{equation}
r \cdot (d_{\text{in}} + d_{\text{out}}) < d_{\text{in}} \cdot d_{\text{out}} \iff r < \frac{d_{\text{in}} \cdot d_{\text{out}}}{d_{\text{in}} + d_{\text{out}}}
\end{equation}

خطای نسبی بازسازی فروبنیوس در این تقریب بر اساس مجموع مقادیر تکین هرس‌شده به دقت قابل محاسبه و کنترل است:
\begin{equation}
\epsilon_{\text{rel}} = \frac{\|\mathbf{W} - \mathbf{U}_{\text{factor}} \mathbf{V}_{\text{factor}}\|_F}{\|\mathbf{W}\|_F} = \sqrt{\frac{\sum_{i=r+1}^{\min(d_{\text{out}}, d_{\text{in}})} \sigma_i^2}{\sum_{i=1}^{\min(d_{\text{out}}, d_{\text{in}})} \sigma_i^2}}
\end{equation}


\paragraph{تحلیل طیفی و نامساوی‌های پایداری خطا}
برای بررسی میزان دقت تقریب ماتریس وزن یا نگاشت توجه در ماژول \lr{direct\_svd}، تفاضل ماتریس فاکتورگیری‌شده از نگاشت متراکم اصلی بر اساس نامساوی هافدینگ-فون‌نیومان ارزیابی می‌گردد. فرض کنید $\mathbf{W} \in \mathbb{R}^{d_{\text{in}} \times d_{\text{out}}}$ ماتریس وزن هدف باشد که توسط تقریب کم‌مرتبه یا پراکنده با رتبه موثر $r$ جایگزین شده است. کران بالای خطای بازسازی بر مبنای نرم فروبنیوس به صورت زیر تصریح می‌شود:
\begin{equation}
\label{eq:spectral_bound_direct_svd}
\| \mathbf{W} - \mathbf{W}_{r} \|_F = \sqrt{\sum_{i=r+1}^{\min(d_{\text{in}}, d_{\text{out}})} \sigma_i^2(\mathbf{W})} \le \epsilon \| \mathbf{W} \|_F
\end{equation}
که در آن $\sigma_i(\mathbf{W})$ بیانگر $i$-امین مقدار منفرد مرتب‌شده ماتریس به صورت نزولی است. نسبت حفظ انرژی طیفی $\eta_r = \frac{\sum_{i=1}^r \sigma_i^2}{\sum_{i=1}^{d} \sigma_i^2}$ تضمین می‌کند که با انتخاب آستانه $\tau \ge 0.90$، بیش از نود درصد از پراکندگی واریانس داده‌ها حفظ شده و افت سیگنال در انتشار گرادیان رو به عقب به حداقل ممکن کاهش می‌یابد.

\paragraph{تحلیل مجانبی و پیچیدگی اشتراک حافظه سخت‌افزاری}
از منظر پیچیدگی محاسباتی و انتقال داده‌ها در سلسومراتب حافظه (حافظه سراسری \lr{HBM} به حافظه محلی \lr{SRAM})، عملیات ضرب ماتریسی در این مکانیزم نیازمند تخصیص بافر موقت با ابعاد فاکتورگیری‌شده است. برای پردازش دنباله‌ای به طول $N$ با بعد مخفی $D$، بار محاسباتی هسته از فرم درجه دو به رابطه زیر تقلیل می‌یابد:
\begin{equation}
\label{eq:complexity_direct_svd}
\mathcal{C}_{\text{direct\_svd}} = \mathcal{O}\left( B \cdot N \cdot D \cdot r + B \cdot H \cdot N \cdot d_k \right) \ll \mathcal{O}\left( B \cdot H \cdot N^2 \cdot d_k \right)
\end{equation}
این بهبود مرتبه پیچیدگی، با کاهش چشمگیر حجم تبادل تانسورها در گذرگاه‌های ارتباطی کارت گرافیک، مانع از اشباع پهنای باند حافظه (پدیده \lr{Memory-Bound}) در فرآیند استنتاج طولانی‌مدت می‌گردد.
